% ---------------------------------------------------------------------
%  Chapitre 3 — Juger les stratégies (version aperçu, rédigée)
% ---------------------------------------------------------------------
\chapter{Juger les stratégies~: l'escalade de prudence}
\label{chap:juger}

Ce chapitre traite la question suivante~: parmi onze stratégies
d'analyse technique, lesquelles produisent un avantage \edgevar{}
significativement positif~? On applique d'abord le test de référence
sous ses hypothèses standard, puis on met à l'épreuve, une à une, ses
conditions d'application, en s'appuyant sur la structure de dépendance
mesurée au chapitre~\ref{chap:donnees}.

% =====================================================================
\section{Étape 1 --- chaque stratégie bat-elle le hasard ?}
\label{sec:b1-etape1}

\subsection{La question posée à l'outil}

La première question posée à l'outil est donc~: l'avantage moyen affiché
par une stratégie est-il réel, ou n'est-il qu'un accident
d'échantillonnage~? On la formalise en confrontant la moyenne de
l'\edgevar{} à la valeur zéro, ce qui conduit au test de Student
présenté ci-dessous.

\subsection{Les méthodes qui y répondent}

\subsubsection*{Test de Student}

Lorsque l'on obtient un edge moyen positif, on cherche à savoir si la
stratégie constitue réellement une prédiction, ou si ce résultat est dû
à l'aléa de l'échantillon. Il faut ici distinguer deux quantités~:
l'edge moyen théorique \(\mu_{\edgevar}\), c'est-à-dire la vraie valeur,
inconnue, qui gouverne la stratégie, et sa moyenne observée \(\bar e\),
calculée sur notre échantillon de trades. On se demande alors si, à
partir de la valeur observée \(\bar e\), on peut conclure que la vraie
valeur \(\mu_{\edgevar}\) est positive.

Le test de Student à un échantillon répond à cette question en
confrontant \(\mu_{\edgevar}\) à une valeur de référence fixée à zéro,
zéro correspondant à l'absence d'avantage sur le hasard. Sa mise en
\oe uvre suppose deux conditions, à savoir que l'edge vérifie un critère
de normalité et que les trades soient indépendants les uns des autres.
La première sera assouplie par le théorème central limite, la seconde,
plus délicate, sera examinée en section suivante. On oppose ainsi deux
hypothèses~:
\[
  \begin{cases}
    \Hzero : \mu_{\edgevar} = 0 & \text{(le vrai edge est nul~: pas mieux que le hasard),}\\[2pt]
    \Hun : \mu_{\edgevar} > 0 & \text{(le vrai edge est positif~: avantage réel).}
  \end{cases}
\]
Le test est unilatéral, car seul un avantage strictement positif nous
intéresse.

Pour trancher, on définit \(t\) comme le nombre d'erreurs-types
nécessaires pour faire passer l'edge moyen de zéro à celui observé dans
notre échantillon~:
\[
  t = \frac{\bar e}{s/\sqrt{n}} \sim t_{n-1} \quad\text{sous } \Hzero .
\]
Sous \Hzero{}, \(t\) suit une loi de Student à \(n-1\) degrés de
liberté. On évalue alors la p-value, c'est-à-dire la probabilité
d'obtenir un \(t\) au moins aussi grand que celui observé si \Hzero{}
était vraie. Si cette p-value est très faible, le hasard n'est plus une
explication crédible de l'edge moyen positif observé, et l'on conclut
que l'edge théorique est positif~: la stratégie fait mieux que le
hasard.

On a dit que le test de Student demandait la normalité de l'edge. On
vérifie cette condition d'abord avec le test de Shapiro--Wilk, puis à
l'aide du théorème central limite. Le test de Shapiro--Wilk a en effet
peu de chances de valider l'hypothèse~: la taille de l'échantillon
déstabilise ce test, alors qu'elle renforce au contraire le théorème
central limite.

\subsubsection*{Test de normalité de Shapiro--Wilk}

Le test de Shapiro--Wilk est un test de normalité qui compare deux
mesures de dispersion, d'une part les écarts à la moyenne, d'autre part
l'ordre des valeurs. Ces deux mesures ne coïncident que si l'hypothèse
de normalité est vraie. On oppose ainsi deux hypothèses~:
\[
  \begin{cases}
    \Hzero \;(W \simeq 1) : \text{l'edge suit une loi normale,}\\[2pt]
    \Hun \;(W \ll 1) : \text{l'edge ne suit pas une loi normale.}
  \end{cases}
\]
On notera que, contrairement au test de Student, c'est ici \Hzero{} qui
correspond à la propriété recherchée~: une p-value faible conduira donc
à rejeter la normalité. On pose alors la statistique
\(W = b^2/\mathrm{SCE}\) comme rapport de ces deux mesures. Lorsque
\(W\) est proche de un, l'échantillon est compatible avec une loi
normale. Pour trancher, on procède comme pour le test de Student~: on
calcule la probabilité d'obtenir notre \(W\) si la loi était normale, et
si cette probabilité est trop faible, on rejette l'hypothèse de
normalité.

Appliqué à nos données, le test rejette la normalité pour les dix
stratégies, les valeurs de \(W\) et les p-values étant détaillées en
annexe~(tableau~\ref{tab:ann-shapiro}). La condition du test de Student
n'est donc, en apparence, pas remplie.

\subsubsection*{Théorème central limite}

Le rejet de Shapiro--Wilk porte sur la non-normalité de l'edge
individuel. Mais comme le test de Student ne dépend que de la moyenne de
l'edge, on peut se rabattre sur le théorème central limite, qui garantit
la quasi-normalité de la moyenne sur les grands échantillons, quelle que
soit la loi de départ~:
\[
  \frac{\bar e - \mu_{\edgevar}}{\sigma/\sqrt{n}}
  \;\xrightarrow[n \to \infty]{}\; \mathcal{N}(0,1) .
\]
Cette justification suppose un effectif suffisant, et un outil générique
devrait écarter les stratégies trop peu fournies. Ici, \(n\) est compris
entre 978 et 8\,513 selon la stratégie, ce qui est largement suffisant.
La condition réellement utile au test de Student est donc satisfaite,
malgré le rejet de Shapiro--Wilk. On remarque de plus que la quantité
ci-dessus est exactement celle du test de Student~: sous \Hzero{}, en
remplaçant l'écart-type théorique \(\sigma\), inconnu, par son
estimation \(s\), on retrouve la statistique
\(t = \bar e/(s/\sqrt{n})\). Le théorème central limite justifie ainsi
directement la forme et la validité du test.

\subsubsection*{Correction de Bonferroni}

Pour appliquer le test de Student, il reste à fixer un seuil pour la
p-value. Pour un test isolé, on prend habituellement \(0{,}05\)~: on
estime que s'il n'y avait que cinq chances sur cent d'obtenir un tel
résultat sous \Hzero{}, il est peu vraisemblable d'être tombé dessus par
hasard, et l'on accepte alors le risque de rejeter \Hzero{}.

Seulement, nous effectuons ici onze tests. Plaçons-nous dans le cas le
plus défavorable, où aucune stratégie n'a d'avantage réel. Pour un test
isolé, la probabilité de ne pas conclure à tort à un avantage est de
\(0{,}95\). En supposant les tests indépendants, la probabilité de ne
commettre aucune erreur sur les onze tests tombe à
\(0{,}95^{11} \approx 0{,}57\). Le risque de commettre au moins un faux
positif sur l'ensemble s'élève donc à \(1 - 0{,}95^{11} \approx
0{,}43\), ce qui est inacceptable.

Pour éviter cela, on raisonne en sens inverse~: on fixe le risque global
à \(0{,}05\) et l'on en déduit le seuil de chaque test. On s'appuie sur
l'inégalité de Boole, selon laquelle la probabilité qu'au moins un
événement se produise est inférieure ou égale à la somme de leurs
probabilités individuelles. Pour garantir un risque global inférieur à
\(0{,}05\), il suffit donc que la somme des seuils individuels vaille
\(0{,}05\). En les prenant tous égaux, on obtient la correction de
Bonferroni~:
\[
  \alpha' = \frac{\alpha}{m} = \frac{0{,}05}{11} \approx 0{,}0045 .
\]

\begin{conditions}
Indépendance des trades (mise à l'épreuve en
section~\ref{sec:b1-dependance})~; normalité de la moyenne assurée par
le théorème central limite ($n$ compris entre 978 et 8\,513 selon la
stratégie).
\end{conditions}

\subsection{Résultats}

\begin{table}[H]
\centering
\caption{Test de l'avantage moyen par stratégie (seuil Bonferroni $\approx 0{,}0045$).}
\label{tab:b1-etape1}
\small
\begin{tabular}{@{}lrrrc@{}}
\toprule
Stratégie & $n$ & \edgevar{} moyen & $p$ (Student) & Verdict \\
\midrule
\texttt{oracle} (témoin) & 2\,152 & $+0{,}2626$ & $5\times10^{-53}$    & bat le hasard \\
\midrule
\texttt{rsi\_classic} & 7\,117 & $+0{,}0123$ & $1{,}2\times10^{-8}$  & bat le hasard \\
\texttt{rsi\_strict}  &   978  & $+0{,}1313$ & $5{,}5\times10^{-10}$ & bat le hasard \\
\texttt{rsi\_trend}   & 2\,789 & $+0{,}0021$ & $0{,}25$             & non \\
\texttt{db\_bottom}   & 8\,513 & $-0{,}0025$ & $0{,}98$             & non \\
\texttt{ma\_crossover}& 5\,417 & $-0{,}0248$ & $\approx 1$          & non \\
\multicolumn{5}{c}{\footnotesize (autres stratégies~: $p>0{,}45$, non significatives)} \\
\bottomrule
\end{tabular}
\end{table}

Outre l'oracle témoin, qui bat le hasard de très loin comme attendu,
seules \texttt{rsi\_classic} et \texttt{rsi\_strict} survivent à la
correction de Bonferroni.

Il est utile de rappeler ce que ce test affirme au juste. Il répond
uniquement à la question de savoir si l'edge moyen positif est dû au
seul échantillonnage, auquel cas on conserve \Hzero{} et la stratégie
n'est pas déclarée gagnante. Dans le cas inverse, on pourra affirmer que
la stratégie a capturé un signal, mais pas qu'elle est gagnante, et cela
pour deux raisons. D'une part, un edge positif ne signifie pas des gains
positifs, mais simplement que l'on perd moins que le hasard. D'autre
part, comme \(n\) est grand, la loi normale de la moyenne se resserre,
et l'on fait ressortir comme significatives des stratégies dont la
moyenne d'edge n'est que très faiblement positive, que les frais
suffiraient à rendre perdantes.

\begin{limites}
À très grand $n$, un test de Student devient sensible au moindre écart~:
une p-value très faible signale un edge statistiquement non nul, sans
garantir qu'il soit économiquement exploitable (frais, exécution).
L'ampleur de l'edge, et non sa seule significativité, devra donc être
prise en compte.
\end{limites}

Ce verdict reste toutefois suspendu à une condition que nous n'avons pas
encore vérifiée, l'indépendance des trades. C'est l'objet de la section
suivante.

% =====================================================================
\section{Étape 1 bis --- le verdict survit-il à la dépendance des trades ?}
\label{sec:b1-dependance}

\subsection{La question posée à l'outil}

On a appliqué le test de Student en supposant les données indépendantes.
Or, comme nous l'avons mesuré au chapitre~\ref{chap:donnees}, nos
données sont dépendantes par le titre, par le marché et d'un mois à
l'autre. Ces dépendances empruntent en réalité deux chemins seulement,
que nous appellerons les deux portes. La première est la porte de
l'action, celle des trades ouverts sur un même titre. La seconde est la
porte de la période, celle des trades ouverts à une même époque, qu'ils
subissent ensemble le mouvement du marché ou qu'ils se prolongent sur
les mêmes mois. Ce sont les deux seules voies par lesquelles deux trades
peuvent partager une cause commune.

Le problème est que le test de Student traite les trades d'une stratégie
comme autant d'informations distinctes, alors que ces informations se
répètent en partie. L'erreur-type s'en trouve sous-estimée. Comme cette
erreur-type sert à mesurer l'écart entre l'avantage observé et la valeur
de référence, un dénominateur trop petit gonfle la statistique de
Student et rend la p-value trop optimiste. Le risque de conclure à tort
qu'une stratégie bat le hasard augmente d'autant.

Nous cherchons donc à savoir si le verdict positif obtenu pour
\texttt{rsi\_classic} et \texttt{rsi\_strict} résiste lorsque l'on
abandonne l'hypothèse d'indépendance. Pour cela, nous corrigeons d'abord
l'erreur-type au moyen de l'effet de plan, puis nous vérifions ce
résultat par un rééchantillonnage qui ne suppose aucune forme
particulière des données. Nous fermons ensuite chaque porte en
regroupant les trades pour que chaque titre, puis chaque période, ne
compte qu'une fois. Nous vérifions enfin que la porte de la période
tient elle-même, en traitant la suite des avantages mensuels comme une
série temporelle.

\subsection{Les méthodes qui y répondent}

\subsubsection*{Corriger l'erreur-type par l'effet de plan}

Nos trades ne forment pas un échantillon tiré au hasard parmi des
observations indépendantes. Ils proviennent d'un ensemble de titres,
chacun livrant plusieurs trades. La théorie des sondages appelle cela un
tirage par grappes, où chaque grappe est ici un titre. Le
chapitre~\ref{chap:donnees} a déjà mesuré, avec la corrélation
intra-titre $\hat\rho$, à quel point les trades d'une même grappe se
ressemblent. Il reste à en tenir compte dans le test.

L'effet de plan, introduit par Kish, chiffre exactement ce que cette
mise en grappes coûte à la précision. Il se calcule ainsi~:
\[
  \mathrm{DEFF} = 1 + (\bar m - 1)\,\hat\rho,
  \qquad n_{\mathrm{eff}} = \frac{n}{\mathrm{DEFF}},
  \qquad \widehat{\mathrm{se}}_{\mathrm{corr}}
        = \frac{s}{\sqrt{n}}\sqrt{\mathrm{DEFF}},
\]
où $\bar m$ désigne le nombre moyen de trades par titre. La variance
réelle de la moyenne vaut $\mathrm{DEFF}$ fois celle que l'on
obtiendrait sous indépendance. Tout se passe donc comme si l'on ne
disposait que de $n_{\mathrm{eff}}$ trades réellement informatifs au
lieu de $n$. Lorsque la corrélation est nulle, l'effet de plan vaut un
et rien ne change. Lorsqu'elle augmente, l'erreur-type est gonflée et le
test devient plus exigeant. On reprend alors le test de Student de
l'étape~1 avec cette erreur-type corrigée, en conservant le même seuil
de Bonferroni.

\subsubsection*{Vérifier sans hypothèse de forme}

Le test précédent s'appuie encore sur la loi de Student, donc sur une
forme particulière des données. Pour nous en affranchir, nous confirmons
le résultat par un rééchantillonnage, la méthode du bootstrap proposée
par Efron~\cite{efron1979}. Le principe consiste à retirer au sort, un
grand nombre de fois, un nouvel échantillon à partir de celui que l'on
possède. Ici nous tirons des titres entiers, avec remise, et jamais des
trades isolés. La structure de dépendance interne de chaque titre est
ainsi transportée telle quelle dans chaque tirage. Après quatre mille
répétitions, la proportion de tirages où l'avantage disparaît fournit
une p-value qui ne suppose ni normalité, ni indépendance entre les
trades d'un même titre. La seule hypothèse qui subsiste est
l'indépendance entre titres différents.

\subsubsection*{Fermer chaque porte en donnant un vote à chaque unité}

Corriger l'erreur-type ne suffit pas. L'effet de plan conserve en effet
la pondération du test initial, dans lequel un titre porteur de deux
cents trades pèse deux cents fois plus qu'un titre n'en portant qu'un
seul. Un avantage concentré sur quelques titres très actifs peut ainsi
passer le test alors qu'il ne se retrouve pas ailleurs. Pour écarter ce
défaut, nous changeons de logique et regroupons les trades avant de
tester, de sorte que chaque unité ne compte qu'une fois.

La première porte, celle de l'action, consiste à calculer d'abord
l'avantage moyen de chaque titre, puis à appliquer le test de Student à
ces quelque cinq cents moyennes. La seconde porte, celle de la période,
calcule de la même manière l'avantage moyen de chaque mois d'entrée,
puis teste ces quelque cent quatre-vingt-dix moyennes. Dans les deux
cas, la dépendance à l'intérieur d'une grappe ne peut plus rien fausser,
puisqu'un titre, ou un mois, ne pèse qu'une seule voix. Le théorème
central limite s'applique alors à ces moyennes de grappes, dont le
nombre reste élevé. Pour être validée, une stratégie doit franchir les
deux portes, chacune au seuil de Bonferroni.

\subsubsection*{Vérifier que la porte de la période tient elle-même}

La porte de la période repose à son tour sur une hypothèse, celle de
mois indépendants. Or nous avons vu au chapitre~\ref{chap:donnees} que
les trades se prolongent souvent au-delà d'un mois, si bien que deux
mois voisins partagent des positions encore ouvertes et se ressemblent.
La suite des avantages mensuels est donc autocorrélée, et le test
précédent pourrait rester trop optimiste. Pour lever ce dernier doute,
nous traitons cette suite comme une série temporelle, en suivant le
déroulement classique de ce type d'analyse.

Avant tout, la série doit être stationnaire, c'est-à-dire garder les
mêmes caractéristiques dans le temps. Nous le vérifions par deux tests
dont les hypothèses de départ sont opposées, celui de Dickey et Fuller
et celui dit KPSS. Lorsque les deux concordent, la conclusion ne dépend
pas du test choisi. Nous examinons ensuite l'autocorrélation de la série
pour déterminer combien de mois passés influencent le mois courant, ce
qui fixe l'ordre d'un modèle autorégressif que nous estimons.

Ce modèle fournit alors une variance de la moyenne qui tient compte de
la persistance, appelée variance de long terme~:
\[
  \gamma_{\mathrm{lr}}
     = \frac{\sigma^2_\eta}{\bigl(1-\sum_{i=1}^{p}\Phi_i\bigr)^{2}},
  \qquad
  \widehat{\mathrm{Var}}(\bar x) \approx \frac{\gamma_{\mathrm{lr}}}{T},
  \qquad
  \mathrm{DEFF}_t = \frac{\gamma_{\mathrm{lr}}}{\gamma(0)} .
\]
Cette variance additionne toutes les autocovariances de la série, là où
la formule usuelle en $s^2/T$ suppose au contraire des mois sans lien.
Plus la persistance est forte, plus la correction est importante. Nous
reprenons enfin le test de Student avec cette variance corrigée, puis
nous vérifions, par le test de Box et Pierce sur les résidus du modèle,
que celui-ci a bien capté toute l'autocorrélation. Si les résidus se
comportent comme un bruit sans mémoire, la correction est justifiée.

\begin{conditions}
Indépendance entre grappes (titres, respectivement mois, ce dernier
point étant précisément l'objet de la vérification par série
temporelle)~; normalité des moyennes de grappes assurée par le théorème
central limite ($k \approx 500$ titres, $T \approx 190$ mois)~;
stationnarité des séries mensuelles vérifiée par Dickey--Fuller et
KPSS~; adéquation des modèles autorégressifs validée par Box--Pierce sur
les résidus.
\end{conditions}

\subsection{Résultats}

\begin{table}[H]
\centering
\caption{Diagnostic de dépendance intra-titre et test corrigé (étape 1b).}
\label{tab:b1-dep-deff}
\small
\begin{tabular}{@{}lrrrccc@{}}
\toprule
Stratégie & $\hat\rho$ & DEFF & $n_{\text{eff}}$ & $p$ naïve & $p$ corrigée & $p$ bootstrap \\
\midrule
\texttt{oracle} (témoin) & $0{,}015$ & $1{,}05$ & $2\,049$ & $\approx 0$ & $\approx 0$ & $<2{,}5\times10^{-4}$ \\
\midrule
\texttt{rsi\_classic} & $0{,}000$ & $1{,}00$ & $7\,117$ & $1{,}2\times10^{-8}$ & $1{,}2\times10^{-8}$ & $<2{,}5\times10^{-4}$ \\
\texttt{rsi\_strict}  & $0{,}134$ & $1{,}15$ & $848$    & $5{,}5\times10^{-10}$ & $7\times10^{-9}$     & $<2{,}5\times10^{-4}$ \\
\multicolumn{7}{c}{\footnotesize (autres stratégies~: verdict « non » inchangé, $\hat\rho \le 0{,}09$)} \\
\bottomrule
\end{tabular}
\end{table}

La corrélation intra-titre étant quasi nulle pour \texttt{rsi\_classic}
et modérée pour \texttt{rsi\_strict}, la correction par l'effet de plan
ne change presque rien~: les deux stratégies conservent leur verdict
positif, et le bootstrap par grappes confirme ces p-values sans recourir
à la loi de Student. À ce stade, la porte de l'action semble donc
franchie. Ce test reste cependant pondéré par le nombre de trades, et ne
dit pas si l'avantage est partagé par l'ensemble des titres ou concentré
sur quelques titres très actifs. C'est ce que le test suivant va
trancher.

\begin{table}[H]
\centering
\caption{Le test des deux portes~: une unité réelle, un vote (étape 1c).}
\label{tab:b1-dep-portes}
\small
\begin{tabular}{@{}lrrrrc@{}}
\toprule
 & \multicolumn{2}{c}{Porte action ($k$ titres)} & \multicolumn{2}{c}{Porte période ($k$ mois)} & \\
\cmidrule(lr){2-3}\cmidrule(lr){4-5}
Stratégie & $k$ & $p_A$ & $k$ & $p_B$ & Verdict \\
\midrule
\texttt{oracle} (témoin) & 482 & $\approx 0$          & 188 & $\approx 0$ & OUI \\
\midrule
\texttt{rsi\_classic}  & 501 & $0{,}059$            & 195 & $0{,}039$ & non \\
\texttt{rsi\_strict}   & 456 & $3{,}9\times10^{-4}$ & 138 & $0{,}26$  & non \\
\texttt{rsi\_trend}    & 495 & $0{,}22$             & 175 & $0{,}29$  & non \\
\texttt{sr\_breakdown} & 457 & $0{,}37$             & 184 & $0{,}94$  & non \\
\multicolumn{6}{c}{\footnotesize (six autres stratégies~: $p_A \ge 0{,}74$, toutes « non »)} \\
\bottomrule
\end{tabular}
\end{table}

La lecture de ce tableau constitue le résultat central du mémoire.
L'oracle, notre témoin, franchit les deux portes de très loin, avec des
p-values proches de zéro. L'outil sait donc dire oui lorsque l'avantage
est réel~: son éventuelle sévérité ne condamne pas d'avance toutes les
stratégies.

Le sort de \texttt{rsi\_classic} est tout autre. La stratégie tombe aux
deux portes, avec des p-values de $0{,}059$ et $0{,}039$, loin du seuil
de $0{,}0045$. Le \(t\) de l'étape~1 valait pourtant $5{,}6$, mais ses
7\,117 trades ne portaient pas 7\,117 informations indépendantes. Dès
que chaque titre ne vote qu'une fois, l'avantage disparaît~: il était
concentré sur les titres et les mois les plus tradés. Le verdict positif
de l'étape~1 était donc un faux positif, que l'outil vient de démasquer
en vérifiant ses propres conditions.

Le cas de \texttt{rsi\_strict} est plus subtil. La stratégie franchit
brillamment la porte de l'action, avec une p-value de
$3{,}9\times10^{-4}$~: son avantage est partagé par de nombreux titres.
Mais elle échoue à la porte de la période, avec une p-value de
$0{,}26$~: présente sur 138~mois seulement, elle concentre son avantage
sur quelques épisodes. Faute de régularité temporelle démontrable, son
signal n'est pas exploitable tel quel.

Au bilan, en dehors du témoin, plus aucune stratégie ne franchit les
deux portes. L'outil discrimine~: il valide l'avantage construit pour
être réel et rejette ceux qui ne tiennent que par la dépendance des
données.

\begin{table}[H]
\centering
\caption{Série mensuelle des edges~: autocorrélation et test corrigé (étape 1d).}
\label{tab:b1-dep-mensuel}
\small
\begin{tabular}{@{}lrrrrrc@{}}
\toprule
Stratégie & $T$ & $\hat\rho(1)$ & $\hat p$ (AR) & $\mathrm{DEFF}_t$ & $p$ corrigée & $p$ Box--Pierce \\
\midrule
\texttt{oracle} (témoin) & 188 & $0{,}03$ & 3 & $2{,}0$ & $2\times10^{-90}$ & $0{,}23$ \\
\midrule
\texttt{ma\_crossover} & 188 & $0{,}56$ & 3 & $6{,}4$ & $0{,}94$ & $0{,}14$ \\
\texttt{dt\_top}       & 197 & $0{,}39$ & 1 & $2{,}2$ & $0{,}97$ & $0{,}69$ \\
\texttt{db\_bottom}    & 197 & $0{,}36$ & 1 & $2{,}1$ & $0{,}82$ & $0{,}90$ \\
\texttt{hs\_inverse}   & 186 & $0{,}31$ & 1 & $1{,}9$ & $0{,}91$ & $0{,}93$ \\
\texttt{rsi\_classic}  & 195 & $0{,}23$ & 1 & $1{,}6$ & $0{,}08$ & $0{,}93$ \\
\texttt{rsi\_strict}   & 138 & $0{,}03$ & 2 & $1{,}5$ & $0{,}30$ & $0{,}96$ \\
\multicolumn{7}{c}{\footnotesize (quatre autres stratégies~: $\mathrm{DEFF}_t \le 2{,}4$, aucune significative)} \\
\bottomrule
\end{tabular}
\end{table}

Les moyennes mensuelles se révèlent bel et bien autocorrélées, jusqu'à
$0{,}56$ pour \texttt{ma\_crossover}~: l'hypothèse de mois indépendants,
sur laquelle reposait la porte de la période, était elle-même optimiste,
la variance étant sous-estimée d'un facteur atteignant $6{,}4$. Toutes
les séries sont stationnaires au vu des tests de Dickey--Fuller et KPSS,
et les modèles autorégressifs sont validés par le test de Box--Pierce
sur leurs résidus~: la correction est donc fondée. Après correction,
aucun verdict ne bascule. \texttt{rsi\_classic} termine à $0{,}08$, loin
du seuil. La porte de la période est fermée proprement. Signalons par
honnêteté que les mois sans trade sont traités comme consécutifs, avec
une couverture de 71~\% pour \texttt{rsi\_strict} et supérieure à 88~\%
ailleurs, limite assumée de cette construction.

\begin{limites}
L'effet de plan de Kish n'est qu'approché pour des grappes de tailles
inégales~; l'agrégation équipondérée protège du faux positif au prix
d'une perte de puissance (choix assumé~: la prudence prime pour un outil
de validation)~; les mois sans trade sont traités comme consécutifs dans
les séries mensuelles~; l'ordre des modèles autorégressifs est plafonné
à~6.
\end{limites}

\subsection{Positionnement dans la littérature financière}
\label{sec:b1-positionnement}

Les outils de cette section ont été construits à partir des seuls
concepts du cours, la théorie des sondages, l'analyse de la variance et
les séries temporelles. Il se trouve qu'ils rejoignent les standards de
l'économétrie financière, signe que la démarche converge vers la
pratique établie. L'agrégation par période, qui calcule une statistique
par mois avant de mener l'inférence sur la suite de ces statistiques,
est exactement la démarche proposée par Fama et
MacBeth~\cite{famamacbeth1973}, devenue un standard du domaine. La
variance de long terme estimée au travers du modèle autorégressif
poursuit le même objectif que les erreurs-types robustes à
l'autocorrélation de Newey et West~\cite{neweywest1987}, dont elle est
une version paramétrique, validée ici par le test de Box--Pierce. Enfin,
le contrôle des faux positifs sur un très grand univers de stratégies
relèverait du \emph{Reality Check} de White~\cite{white2000}~; avec onze
stratégies déclarées à l'avance, la correction de Bonferroni suffit et
demeure le choix le plus prudent.

Le verdict au standard corrigé est donc le suivant~: aucune des dix
stratégies réelles ne démontre d'avantage. Le verdict de l'étape~1 ne
tenait que par une condition non vérifiée. C'est le résultat central de
ce mémoire~: l'outil de validation a démasqué son propre faux positif en
testant ses conditions d'application. Un verdict négatif rendu
proprement est la preuve que l'outil fonctionne.

% =====================================================================
\section{Lecture du verdict --- le marché comme série temporelle}
\label{sec:b1-lecture}

Le verdict est tombé, aucune stratégie ne démontre d'avantage. Avant de
conclure, on peut se demander si ce résultat était prévisible, en
interrogeant le marché lui-même avec les outils des séries temporelles.
Le chapitre~\ref{chap:donnees} a montré que la dépendance simultanée est
massive, le facteur marché portant 73~\% des mouvements. Reste à savoir
s'il existe de la dépendance retardée, la seule qu'une stratégie puisse
exploiter, puisqu'il faut un délai pour passer un ordre.

On vérifie d'abord la stationnarité, condition de toute modélisation,
par le test de Dickey--Fuller augmenté, appliqué au prix puis au
rendement de l'indice de marché. Le prix ne rejette pas la racine
unitaire, avec une p-value de $0{,}999$~: il se comporte comme une
marche aléatoire, dont le niveau de demain est celui d'aujourd'hui plus
un choc imprévisible. Le rendement, lui, est nettement stationnaire,
avec une p-value de l'ordre de $10^{-26}$. On modélise donc le
rendement, jamais le prix. Son autocorrélation est quasi nulle à tous
les retards examinés, entre $-0{,}08$ et $+0{,}09$ sur les dix premiers
jours. Un modèle ARIMA~\cite{boxjenkins1970} ajusté sur ce rendement ne
trouve d'ailleurs rien à saisir, ses deux termes, autorégressif et
moyenne mobile, se compensant presque exactement.

Le rendement du marché est ainsi proche d'un bruit sans mémoire~: son
propre passé ne prédit guère le lendemain. C'est ce que la littérature
appelle l'efficience faible du marché. La boucle est alors bouclée. Le
marché regorge de dépendance simultanée, celle qui fausse les tests et
qu'aucune stratégie ne peut exploiter faute de délai, et manque
précisément de dépendance retardée, la seule qui serait exploitable.
L'échec des dix stratégies n'est pas un accident~: c'est la signature de
cette asymétrie.

\begin{conditions}
Travailler sur les rendements (stationnaires), jamais sur les prix
bruts~; la stationnarité est vérifiée par le test de Dickey--Fuller
avant toute modélisation.
\end{conditions}

% =====================================================================
\section{Synthèse du chapitre}
\label{sec:b1-synthese}

Retraçons le chemin parcouru. Le test de référence déclarait deux
stratégies gagnantes sur dix. La mise à l'épreuve de ses conditions
d'application, par l'effet de plan, le bootstrap par grappes,
l'agrégation en deux portes et la correction de la série mensuelle, a
renversé ces deux verdicts. \texttt{rsi\_classic} était un faux positif
né de la répétition d'une même information sur des milliers de trades.
\texttt{rsi\_strict} possède un avantage réel à l'échelle des titres,
mais trop épisodique dans le temps pour être retenu. La lecture du
marché lui-même rend ce résultat cohérent~: un marché proche de
l'efficience faible n'offre presque pas de dépendance retardée à
exploiter. Le message central est ailleurs que dans les stratégies. Le
protocole a fait exactement son travail, il a démasqué son propre faux
positif en vérifiant ses conditions, tout en validant sans hésiter le
témoin dont l'avantage était réel. Un verdict négatif rendu proprement,
c'est l'outil qui fonctionne.
